\chapter{晶体上的 t-结构}
\markboth{第 4 章\ 晶体上的 t-结构}{第 4 章\ 晶体上的 t-结构}

晶体范畴有两个自然的 t-结构，分别与左实现和右实现相容。右实现的主要优点之一，是与它相容的 t-结构性质好得多。

本节将定义这两个 t-结构，并证明它们的一些基本性质。这些性质包括：归纳/遗忘函子的左/右 t-正合性与上同调幅度的有界性；DG 概形的 $\Crys$ 的左完备性；与 t-结构心的导出范畴的关系。

\section{左 t-结构}

本小节中，我们不作预栈局部几乎有限型的假设。

\ob{4.1.1}{n:4.1.1}回忆 [GL:QCoh, Sec.~1.2.3]：对任意预栈 $Z$，范畴 $\QCoh(Z)$ 有典范 t-结构，其特征如下：对象 $F\in\QCoh(Z)$ 属于 $\QCoh(Z)^{\leq 0}$，当且仅当对每个 $S\in\DGSch_{\mathrm{aff}}$ 及映射 $\varphi:S\to Z$，有
\[
\varphi^*(F)\in\QCoh(S)^{\leq 0}.
\]
特别地，取 $Z=Y_{\dR}$，我们得到 $\Crys^l(Y)$ 上的典范 t-结构，称之为「左 t-结构」。

按定义，函子
\[
\oblv^l:\Crys^l(X)\to\QCoh(X)
\]
对左 t-结构是右 t-正合的。

\ob{4.1.2}{n:4.1.2}一般说来，左 t-结构的性质相当差。不过，我们有如下断言：

\begin{othm}{命题}{4.1.3}{n:4.1.3}
设 $Y$ 是经典形式光滑的预栈。则
\[
M\in\Crys^l(Y)^{\leq 0}\iff \oblv^l_Y(M)\in\QCoh(Y)^{\leq 0}.
\]
\end{othm}

\begin{proof}
需要证明：若 $M\in\Crys^l(Y)$ 满足 $\oblv^l_Y(M)\in\QCoh(Y)^{\leq 0}$，则 $M\in\Crys^l(Y)^{\leq 0}$。也就是说，需要证明对每个 $S\in\DGSch_{\mathrm{aff}}$ 及 $\varphi:S\to Y_{\dR}$，有 $\varphi^*(M)\in\QCoh(S)^{\leq 0}$。

记 $Y_\bullet/Y_{\dR}$ 为映射 $p_{\dR,Y}:Y\to Y_{\dR}$ 的 Čech 神经。由引理 1.2.4，存在映射 $\varphi':S\to Y$ 及同构 $\varphi\simeq p_{\dR,Y}\circ\varphi'$。断言现在由 $\varphi'^*$ 右 t-正合推出。
\end{proof}

\section{右 t-结构}

从此刻到本节末，我们恢复所有预栈局部几乎有限型的假设（除明确说明外）。

本小节专门讨论 DG 概形的情形。

\ob{4.2.1}{n:4.2.1}设 $X$ 是 DG 概形。回忆范畴 $\IndCoh(X)$ 有自然的、与滤过余极限相容的 t-结构，见 [IndCoh, Sect.~1.2]。

它的特征性质是：$\IndCoh(X)$ 中的对象是连接的（即落在 $\IndCoh(X)^{\leq 0}$ 中），当且仅当它在函子 $\Psi_X:\IndCoh(X)\to\QCoh(X)$ 下的像连接。

\ob{4.2.2}{n:4.2.2}我们通过宣布
\[
M\in\Crys^r(X)^{\geq 0}\iff \oblv^r_X(M)\in\IndCoh(X)^{\geq 0}
\]
来定义 $\Crys^r(X)$ 上的右 t-结构。

下文中，我们把 $\Crys^r(X)$ 上的右 t-结构称为晶体上的「那个」t-结构。换言之，默认情况下我们考虑的 t-结构是右 t-结构。按构造，这个 t-结构也与滤过余极限相容，因为 $\oblv^r_X$ 连续。

\begin{annote}
「那个」t-结构这个叫法来自 §0.2.1：在光滑经典概形上，右实现给出的 t-结构在 Riemann--Hilbert 对应下就是右 D-模的反常 t-结构，而左 t-结构只反映拟凝聚层的通常 t-结构。因此凡不加限定时，本文默认使用右 t-结构。\hfill【来源：原书 §0.2.1、§4.2.2（A）】
\end{annote}


\ob{4.2.3}{n:4.2.3}我们断言 $\Crys^r(X)$ 上的右 t-结构是 Zariski 局部的，即一个对象连接/余连接，当且仅当它限制到某个 Zariski 覆盖后具有该性质。事实上，这由 $\IndCoh(X)$ 上 t-结构的相应性质推出，见 [IndCoh, Corollary 4.2.3]。

\ob{4.2.4}{n:4.2.4}右 t-结构与 Kashiwara 引理。设 $i:X\to Z$ 是 DG 概形的闭浸入。记 $i_{\dR,*}$ 为函子 $\Crys^r(X)\to\Crys^r(Z)$，它等于复合
\[
\Crys^r(X)\xrightarrow{(i'^{\dagger,r})^{-1}}\Crys^r(Z)_X\hookrightarrow\Crys^r(Z),
\]
按构造，它是 $i^{\dagger,r}$ 的左伴随。

有：

\begin{othm}{命题}{4.2.5}{n:4.2.5}
函子 $i_{\dR,*}:\Crys^r(X)\to\Crys^r(Z)$ 是 t-正合的。
\end{othm}

\begin{proof}
注意全子范畴
\[
\Crys^r(Z)_X\subset\Crys^r(Z)
\]
与 t-结构相容，因为它是函子 $j^{\dagger,r}$ 的核，而后者 t-正合（这里 $j$ 表示开浸入 $Z-X\hookrightarrow Z$）。

于是只需证明函子
\[
i'^{\dagger,r}:\Crys^r(Z)_X\to\Crys^r(X)
\]
t-正合。

因此需要证明：对 $M\in\Crys^r(Z)_X$，有
\[
\oblv^r_Z(M)\in\IndCoh(Z)^{>0}\iff \oblv^r_X(i^{\dagger,r}(M))\in\IndCoh(X)^{>0}.
\]
回忆记号
\[
\IndCoh(Z)_X:=\ker\big(j^!:\IndCoh(Z)\to\IndCoh(Z-X)\big).
\]
只需证明下述引理：
\end{proof}

\begin{othm}{引理}{4.2.6}{n:4.2.6}
设 $i:X\to Z$ 是闭浸入。则对 $F\in\IndCoh(Z)_X$ 有
\[
F\in\IndCoh(Z)^{>0}\iff i^!(F)\in\IndCoh(X)^{>0}.
\]
\end{othm}

\begin{proof}[引理 4.2.6 的证明]
「$\Rightarrow$」由 $i^!$ 左 t-正合推出（它是 t-正合函子 $i_*^{\IndCoh}$ 的右伴随）。

对反向蕴含，注意全子范畴
\[
\IndCoh(Z)_X\subset\IndCoh(Z)
\]
也与 t-结构相容，因为它是 t-正合函子 $j^!$ 的核。此外，由 [IndCoh, Proposition 4.1.7(b)]，$\IndCoh(Z)_X$ 上的 t-结构由
\[
\Coh(Z)_X:=\ker\big(j^*:\Coh(Z)\to\Coh(Z-X)\big)
\]
上的 t-结构生成。

设 $F\in\IndCoh(Z)_X$ 满足 $i^!(F)\in\IndCoh(X)^{>0}$。需要证明 $F$ 与 $(\Coh(Z)_X)^{\leq 0}$ 右正交。按假设，$F$ 与 $\Coh(X)^{\leq 0}$ 在
\[
\Coh(X)\xrightarrow{i_*}\Coh(Z)_X\to\IndCoh(Z)_X
\]
下的本质像右正交。然而容易看出，$(\Coh(Z)_X)^{\leq 0}$ 的每个对象都可以由 $\Coh(X)^{\leq 0}$ 的本质像中的对象经过有限次逐步扩张得到，这蕴含所需断言。
\end{proof}

\begin{othm}{推论}{4.2.7}{n:4.2.7}
若 DG 概形的映射 $X_1\to X_2$ 诱导同构
\[
\clred X_1\to\clred X_2,
\]
则 $\Crys^r(X_1)\simeq\Crys^r(X_2)$ 上相应的 t-结构重合。
\end{othm}

\ob{4.2.8}{n:4.2.8}设 $X$ 是 DG 概形。按构造，遗忘函子 $\oblv^r_X$ 左 t-正合。于是由伴随性，函子 $\indf^r_X$ 右 t-正合。

现在断言：

\begin{othm}{命题}{4.2.9}{n:4.2.9}
函子 $\indf^r_X$ 是 t-正合的。
\end{othm}

\begin{proof}
只需证明复合 $\oblv^r_X\circ\indf^r_X$ 左 t-正合。我们由命题 3.3.2(b) 推出这一点：

函子 $p_s^!$ 左 t-正合（例如对 $\Delta_X:X\to X\times X$ 应用引理 4.2.6）。函子 $(p_t)_*^{\IndCoh}$ 由 [GR1, Lemma 2.7.11] 左 t-正合（事实上 t-正合）。
\end{proof}

\ob{4.2.10}{n:4.2.10}现在断言：

\begin{othm}{命题}{4.2.11}{n:4.2.11}
(a) 若 $X$ 是光滑经典概形，则 $\oblv^r_X$ t-正合。

(b) 对拟紧 DG 概形 $X$，函子 $\oblv^r_X$ 有界上同调幅度。
\end{othm}

\begin{proof}
设 $X$ 是光滑经典概形。按 $\Crys^r(X)$ 上 t-结构的定义，$\IndCoh(X)^{\leq 0}$ 在 $\indf^r_X$ 下的本质像通过取余极限生成 $\Crys^r(X)^{\leq 0}$。因此，为证明 $\oblv^r_X$ 右 t-正合，只需对函子 $\oblv^r_X\circ\indf^r_X$ 证明同样的事情。我们将由命题 3.3.2(b) 推出这一点：

可以写
\[
(X\times X)^{\wedge}_X\simeq\varinjlim_n X_n,
\]
其中 $X_n\to X\times X$ 是对角线的第 $n$ 个无穷小邻域。于是由 [GR1, Equation (2.2)]，
\[
(p_t)_*^{\IndCoh}\circ p_s^!\simeq\varinjlim_n (p_t\circ i_n)_*^{\IndCoh}\circ(p_s\circ i_n)^!.
\]
现在，每个函子 $(p_t\circ i_n)_*^{\IndCoh}$ 由 [GR1, Lemma 2.7.11] t-正合；每个函子 $(p_s\circ i_n)^!$ t-正合，因为 $p_s\circ i_n:X_n\to X$ 有限平坦。

现在设 $X$ 是拟紧 DG 概形，来证明 $\oblv^r_X$ 有界上同调幅度。问题容易化归到 $X$ 仿射的情形；取闭浸入 $i:X\hookrightarrow Z$，其中 $Z$ 光滑。由命题 4.2.5 与 (a)，只需证明函子 $i^!:\IndCoh(Z)\to\IndCoh(X)$ 有界上同调幅度，而后者容易由 $Z$ 正则推出。
\end{proof}

\section{ind-概形上晶体的右 t-结构}

\ob{4.3.1}{n:4.3.1}设 $X$ 是 DG ind-概形。固定 $X$ 的一个表示
\begin{equation}
X=\varinjlim_\alpha X_\alpha\tag{4.1}
\end{equation}
如 [GR1, Prop.~1.7.6]。对每个 $\alpha$，记 $i_\alpha$ 为相应的闭浸入 $X_\alpha\to X$；对每个 $\alpha_1\to\alpha_2$，记 $i_{\alpha_1,\alpha_2}$ 为闭浸入 $X_{\alpha_1}\to X_{\alpha_2}$。

有：
\[
\Crys^r(X)\simeq\varprojlim_\alpha\Crys^r(X_\alpha),
\]
其中对 $\alpha_1\to\alpha_2$，函子 $\Crys^r(X_{\alpha_2})\to\Crys^r(X_{\alpha_1})$ 由 $i^{\dagger,r}_{\alpha_1,\alpha_2}$ 给出。

于是由 [GL:DG, Sect.~1.3.3]，
\[
\Crys^r(X)\simeq\varinjlim_\alpha\Crys^r(X_\alpha),
\]
其中对 $\alpha_1\to\alpha_2$，函子 $\Crys^r(X_{\alpha_1})\to\Crys^r(X_{\alpha_2})$ 由 $(i_{\alpha_1,\alpha_2})_{\dR,*}$ 给出。

特别地，对每个 $\alpha$ 得到一对伴随函子
\[
(i_\alpha)_{\dR,*}:\Crys^r(X_\alpha)\rightleftarrows\Crys^r(X):i^{\dagger,r}_\alpha.
\]

\ob{4.3.2}{n:4.3.2}回忆 [GR1, Sect.~2.5]：$\IndCoh(X)$ 有自然的、与滤过余极限相容的 t-结构。

利用 $\IndCoh(X)$ 上的这个 t-结构，可以定义 $\Crys^r(X)$ 上的右 t-结构。即
\[
M\in\Crys^r(X)^{\geq 0}\iff \oblv^r_X(M)\in\IndCoh(X)^{\geq 0}.
\]
由于 $\oblv^r_X$ 保持余极限，这个 t-结构与滤过余极限相容。利用表示 (4.1)，可以更具体地描述这个 t-结构，方式类似于 [GR1, Lemma 2.5.3] 中对 $\IndCoh(X)$ 上 t-结构的描述：

\begin{othm}{引理}{4.3.3}{n:4.3.3}
在上述情形下有：

(a) 对象 $F\in\Crys^r(X)$ 属于 $\Crys^r(X)^{\geq 0}$，当且仅当对每个 $\alpha$，对象 $i^{\dagger,r}_\alpha(F)\in\Crys^r(X_\alpha)$ 属于 $\Crys^r(X_\alpha)^{\geq 0}$。

(b) 范畴 $\Crys^r(X)^{\leq 0}$ 由函子 $(i_\alpha)_{\dR,*}\big(\Crys^r(X_\alpha)^{\leq 0}\big)$ 的本质像在余极限下生成。
\end{othm}

\begin{proof}
(a) 由定义与 [GR1, Lemma 2.5.3(a)] 推出。(b) 由 (a) 形式地推出。
\end{proof}

\ob{4.3.4}{n:4.3.4}设 $i:X\to X$ 是把 DG 概形闭浸入 DG ind-概形。由与 [GR1, Lemma 2.5.5] 完全相同的论证，有：

\begin{othm}{引理}{4.3.5}{n:4.3.5}
函子 $i_{\dR,*}$ 是 t-正合的。
\end{othm}
\ob{4.3.6}{n:4.3.6}作为上述 DG ind-概形上右晶体 t-结构行为的一个说明，考虑如下情形。设 $i:X\to Z$ 是拟紧 DG 概形的闭浸入。记 $Y$ 为 $X$ 在 $Z$ 中的形式完备化，看作 $\DGIndSch_{\mathrm{laft}}$ 的对象；记 $i':X\to Y$ 为相应的映射。

我们断言：

\begin{othm}{引理}{4.3.7}{n:4.3.7}
由同构 $i'_{\dR}:X_{\dR}\to Y_{\dR}$ 诱导的等价 $\Crys^r(X)\simeq\Crys^r(Y)$ 与 t-结构相容。
\end{othm}

\begin{proof}[证明一]
由命题 4.2.5 以及 [GR1, Proposition 7.4.5] 中的等价
\[
\IndCoh(Y)\simeq\IndCoh(Z)_Y
\]
与 t-结构相容（见 [GR1, Lemma 7.4.8]）推出。
\end{proof}

\begin{proof}[证明二]
由交换图
\[
\begin{tikzcd}[column sep=large, row sep=large]
 \Crys^r(X) & \Crys^r(Y) \arrow[l, "i'^{\dagger,r}"'] \\
 \IndCoh(X) \arrow[u, "\oblv^r_X"] & \IndCoh(Y) \arrow[l, "i'^!"] \arrow[u, "\oblv^r_Y"']
\end{tikzcd}
\]
得到：只需证明对 $F\in\IndCoh(Y)$ 有
\[
F\in\IndCoh(Y)^{>0}\iff i'^!(F)\in\IndCoh(X)^{>0},
\]
这由 [GR1, Lemma 7.4.8] 与引理 4.2.6 形式地推出（也可以直接证明）。
\end{proof}

\section{左 t-结构的进一步性质}

\ob{4.4.1}{n:4.4.1}首先描述光滑经典概形情形下晶体上左、右 t-结构之间的关系。

\begin{othm}{命题}{4.4.2}{n:4.4.2}
设 $X$ 是 $n$ 维光滑经典概形。则
\[
F\in\Crys^l(X)^{\leq 0}\iff F\in\Crys^r(X)^{\leq -n}.
\]
也就是说，左 t-结构与右 t-结构相差 $X$ 的维数这样一个平移。
\end{othm}

\begin{proof}
回忆两个遗忘函子由交换图
\[
\begin{tikzcd}[column sep=large, row sep=large]
 & \Crys(X) \arrow[dl, "\oblv^l_X"'] \arrow[dr, "\oblv^r_X"] & \\
 \QCoh(X) \arrow[rr, "\Upsilon_X"'] & & \IndCoh(X)
\end{tikzcd}
\]
相联系。

当 $X$ 是 $n$ 维光滑经典概形时，函子 $\Upsilon_X$ 是等价，并把 $\QCoh(X)^{\leq 0}$ 同构地送到 $\IndCoh(X)^{\leq -n}$。断言现在由命题 4.1.3 与命题 4.2.11(a)，结合 $\oblv^r_X$ 保守这一事实推出。
\end{proof}

\ob{4.4.3}{n:4.4.3}下一个命题对任意 DG 概形 $X$ 比较 $\Crys(X)$ 上的「左」t-结构与「右」t-结构。

\begin{othm}{命题}{4.4.4}{n:4.4.4}
设 $X$ 拟紧。则恒等函子
\[
\Crys^l(X)\to\Crys^r(X)
\]
有界幅度，即左、右 t-结构之间的差别有界。
\end{othm}

\begin{proof}
不妨设 $X$ 仿射。设 $Z$ 是 $n$ 维光滑经典概形，$i:X\hookrightarrow Z$ 是闭浸入。我们断言：对 $M^l\in\Crys^l(X)$ 及相应的对象 $M^r\in\Crys^r(X)$，有
\begin{equation}
M^l\in(\Crys^l(X))^{\leq 0}\Rightarrow M^r\in(\Crys^r(X))^{\leq 0},\qquad
M^r\in(\Crys^r(X))^{\leq 0}\Rightarrow M^l\in(\Crys^l(X))^{\leq n}.\tag{4.2}
\end{equation}
记 $j:U\hookrightarrow Z$ 为补开浸入。记 $Y$ 为 $X$ 在 $Z$ 中的形式完备化；记 $\widehat{i}:Y\to Z$。

映射 $X\to Y$ 定义同构 $X_{\dR}\to Y_{\dR}$，从而可以认同 $\Crys^l(X)\simeq\Crys^l(Y)$。应用命题 4.1.3，得到
\begin{equation}
M^l\in(\Crys^l(X))^{\leq 0}\iff \oblv^l_Y(M^l)\in\QCoh(Y)^{\leq 0},\tag{4.3}
\end{equation}
其中 $\QCoh(Y)$ 上的 t-结构是 §4.1.1 中的那个。

考虑子范畴 $\QCoh(Z)_X\subset\QCoh(Z)$，按定义它等于
\[
\ker\big(j^*:\QCoh(Z)\to\QCoh(U)\big).
\]
这个子范畴与 $\QCoh(Z)$ 上的 t-结构相容，因为函子 $j^*$ t-正合。

回忆（见 [GR1, Proposition 7.1.3]）：函子 $\widehat{i}_*$ 定义等价
\[
\QCoh(Z)_X\to\QCoh(Y).
\]
设 $F$ 是 $\QCoh(Z)_X$ 中对应于 $\oblv^l_Y(M^l)\in\QCoh(Y)$ 的对象。有
\[
\Upsilon_Z(F)\simeq\oblv^r_Z(i_{\dR,*}(M^r)).
\]
由于函子 $i_{\dR,*}$ t-正合（命题 4.2.5），且 $\Upsilon_Z$ 把上同调次数平移 $[-n]$，得到
\begin{equation}
M^r\in(\Crys^r(X))^{\leq 0}\iff F\in(\QCoh(Z)_X)^{\leq n}.\tag{4.4}
\end{equation}
结合 (4.3) 与 (4.4)，(4.2) 中的蕴含由下述断言推出：
\end{proof}

\begin{othm}{引理}{4.4.5}{n:4.4.5}
等价 $\widehat{i}_*:\QCoh(Z)_X\simeq\QCoh(Y)$ 关于 $\QCoh(Z)_X$ 上从 $\QCoh(Z)$ 继承的 t-结构与 §4.1.1 中 $\QCoh(Y)$ 上的 t-结构，具有如下性质：

(a) 若 $F\in(\QCoh(Z)_X)^{\leq 0}$，则 $\widehat{i}_*(F)\in\QCoh(Y)^{\leq 0}$。

(b) 若 $\widehat{i}_*(F)\in\QCoh(Y)^{\leq 0}$，则 $F\in(\QCoh(Z)_X)^{\leq n}$。
\end{othm}

\begin{proof}[引理 4.4.5 的证明]
(a) 由函子 $\widehat{i}^*$ 右 t-正合推出。

为证明 (b)，注意范畴 $\QCoh(Y)^{\leq 0}$ 在取余极限下由 $\QCoh(Z)^{\leq 0}$ 在函子 $\widehat{i}_*$ 下的本质像生成，见 [GR1, Proposition 7.3.5]。因此只需证明函子
\[
\widehat{i}_*:\QCoh(Z)\to\QCoh(Y)\simeq\QCoh(Z)_X
\]
的上同调幅度被 $n$ 界住。然而上述函子是嵌入
\[
\QCoh(Z)_X\hookrightarrow\QCoh(Z)
\]
的右伴随，由
\[
F'\mapsto\Cone\big(F'\to j_*j^*(F')\big)[-1]
\]
给出。现在 $j^*$ t-正合，$j_*$ 的上同调幅度被 $n-1$ 界住。这蕴含所需断言。
\end{proof}

\ob{4.4.6}{n:4.4.6}设 $X$ 是任意拟紧 DG 概形。有：

\begin{othm}{命题}{4.4.7}{n:4.4.7}
(a) 函子 $\oblv^l_X:\Crys(X)\to\QCoh(X)$ 有界上同调幅度。

(b) 若 $X$ 最终余连接，则函子 $\indf^l_X:\QCoh(X)\to\Crys(X)$ 的上同调幅度上有界。
\end{othm}

\begin{proof}
对 (a)，不妨设 $X$ 仿射，并取闭浸入 $i:X\hookrightarrow Z$，其中 $Z$ 是光滑经典概形。此时断言由命题 4.2.5 与函子
\[
i^*:\QCoh(Z)\to\QCoh(X)
\]
有界上同调幅度推出。

(b) 由 (a) 与 $(\indf^l_X,\oblv^l_X)$-伴随性推出。
\end{proof}

\begin{othm}{注}{4.4.8}{n:4.4.8}
(b) 中 $X$ 最终余连接的假设是本质的；否则可以用例 3.4.5 中的 DG 概形给出反例。此外，容易证明：$\indf^l_X$ 的上同调幅度下有界，当且仅当 $X$ 是 Gorenstein 的（见引理 4.6.12）。
\end{othm}

\section{左完备性}

\ob{4.5.1}{n:4.5.1}设 $X$ 是仿射光滑经典概形。我们注意到，此时范畴 $\Crys^r(X)$ 含有一个典范对象
\[
\indf^r_X(\mathcal{O}_X),
\]
它落在 t-结构的心中（见命题 4.2.9），并且是投射的，即
\[
H^0(N)=0\Rightarrow\Hom_{\Crys^r(X)}\big(\indf^r_X(\mathcal{O}_X),N\big)=0.
\]
此外，$\indf^r_X(\mathcal{O}_X)$ 是 $\Crys^r(X)$ 的紧生成元。这蕴含：

\begin{othm}{推论}{4.5.2}{n:4.5.2}
设 $X$ 是仿射光滑经典概形。则范畴 $\Crys^r(X)$ 在其 t-结构下左完备。
\end{othm}

\ob{4.5.3}{n:4.5.3}上述推论对任意 DG 概形 $X$ 蕴含左完备性：

\begin{othm}{推论}{4.5.4}{n:4.5.4}
对任意 DG 概形 $X$，范畴 $\Crys^r(X)$ 在「右」t-结构下左完备。
\end{othm}

\begin{proof}
首先注意，左完备性是 Zariski 局部的（用与 [GL:QCoh, Proposition 5.2.4] 相同的论证证明）。于是不妨设 $X$ 仿射。取闭浸入 $i:X\hookrightarrow Z$，其中 $Z$ 是光滑经典概形。此时断言由以下事实形式地推出：函子 $i_{\dR,*}$ 连续、全忠实（由命题 2.5.6）、t-正合（由命题 4.2.5），且 $\Crys^r(Z)$ 左完备（由上一推论）。

这里还有一个替代论证：

由推论 4.2.7，可以假设 $X$ 最终余连接。此时函子 $\oblv^l_X$ 与极限交换，因为它有左伴随。此外，由引理 2.2.6，$\oblv^l_X$ 保守；由命题 4.4.7，它有界上同调幅度。因此，$\QCoh(X)$ 在其 t-结构下左完备这一事实蕴含 $\Crys^r(X)$ 的相应事实。
\end{proof}

\begin{othm}{注}{4.5.5}{n:4.5.5}
右完备性不成问题：由于我们的 t-结构与滤过余极限相容，右完备性等价于 t-结构在余连接子范畴上可分，而这是显然的，因为 $\oblv^r_X$ 左 t-正合且保守，并且 $\IndCoh(X)^+\simeq\QCoh(X)^+$ 上的 t-结构具有该性质。
\end{othm}

\ob{4.5.6}{n:4.5.6}把推论 4.5.4 与命题 4.4.4 结合，得到：

\begin{othm}{推论}{4.5.7}{n:4.5.7}
对拟紧 DG 概形 $X$，范畴 $\Crys(X)$ 在「左」t-结构下也左完备。
\end{othm}

\section{「粗糙的」归纳函子与遗忘函子}

\ob{4.6.1}{n:4.6.1}设 $X$ 是 DG 概形。回忆函子 $\Psi_X$ 把范畴 $\QCoh(X)$ 与 $\IndCoh(X)$ 的左完备化等同起来（见 [IndCoh, Proposition 1.3.4]）。

由于范畴 $\Crys^r(X)$ 在其 t-结构下左完备，且函子 $\indf^r_X$ t-正合，由左完备化的泛性质得到：

\begin{othm}{推论}{4.6.2}{n:4.6.2}
函子 $\indf^r_X$ 典范地分解为
\[
\IndCoh(X)\xrightarrow{\Psi_X}\QCoh(X)\xrightarrow{{}'\indf^r_X}\Crys^r(X).
\]
\end{othm}

\ob{4.6.3}{n:4.6.3}也可以考虑函子
\[
{}'\oblv^r_X:\Crys^r(X)\to\QCoh(X),
\]
它由 $\Psi_X\circ\oblv^r_X$ 给出，其中 $\Psi_X:\IndCoh(X)\to\QCoh(X)$ 是 [IndCoh, Sect.~1.1.5] 中的函子。

显然函子 ${}'\oblv^r_X$ 有有限上同调幅度。事实上，这由 $\oblv^r_X$ 的相应事实与 $\Psi_X$ t-正合（见 [IndCoh, Lemma 1.2.2]）推出。

\begin{othm}{命题}{4.6.4}{n:4.6.4}
函子 ${}'\oblv^r_X$ 保守。
\end{othm}

\begin{proof}
断言是 Zariski 局部的，于是不妨设 $X$ 仿射。取闭浸入 $i:X\to Z$，其中 $Z$ 是光滑经典仿射概形。

记 $i^{\QCoh,!}:\QCoh(Z)\to\QCoh(X)$ 为 $i_*:\QCoh(X)\to\QCoh(Z)$ 的右伴随\footnote{尽管这与我们无关，我们仍指出 $i^{\QCoh,!}$ 连续。这是因为函子 $i_*:\QCoh(X)\to\QCoh(Z)$ 把紧对象送到紧对象（由于 $Z$ 正则，其上任一凝聚层都是完美的）。}。容易看出有函子的典范同构
\[
\Psi_X\circ i^!\simeq i^{\QCoh,!}\circ\Psi_Z.
\]
于是对 $M\in\Crys^r(Z)$，有
\[
{}'\oblv^r_X(i^{\dagger,r}(M))\simeq i^{\QCoh,!}({}'\oblv^r_Z(M)).
\]
应用 Kashiwara 引理，命题的断言由下述引理推出：
\end{proof}

\begin{othm}{引理}{4.6.5}{n:4.6.5}
函子 $i^{\QCoh,!}:\QCoh(Z)\to\QCoh(X)$ 限制到 $\QCoh(Z)_X$ 上时保守。
\end{othm}

\begin{proof}[引理 4.6.5 的证明]
需要证明函子 $i_*:\QCoh(X)\to\QCoh(Z)$ 的本质像生成 $\QCoh(Z)_X$。

首先断言 $\QCoh(Z)_X$ 由有界对象的子范畴 $(\QCoh(Z)_X)^b$ 生成。这由 $\QCoh(Z)$ 的相应事实以及嵌入 $\QCoh(Z)_X\hookrightarrow\QCoh(Z)$ 有有界上同调幅度的右伴随推出。由 dévissage得到：$\QCoh(Z)_X$ 由 $(\QCoh(Z)_X)^{\heartsuit}$ 生成，进而由 $(\QCoh(Z)_X)^{\heartsuit}\cap\Coh(Z)$ 生成。

然而显然 $(\QCoh(Z)_X)^{\heartsuit}\cap\Coh(Z)$ 的每个对象都是 $\Coh(X)^{\heartsuit}$ 的本质像中对象的有限次扩张。
\end{proof}

\begin{othm}{注}{4.6.6}{n:4.6.6}
当 $X$ 最终余连接时，我们将在下面给出命题 4.6.4 的一个更干净的证明。
\end{othm}

\ob{4.6.7}{n:4.6.7}现在假设 $X$ 最终余连接。回忆此时函子 $\Psi_X$ 有全忠实的左伴随 $\Xi_X$（见 [IndCoh, Proposition 1.5.3]）。

我们注意到：

\begin{othm}{引理}{4.6.8}{n:4.6.8}
存在典范同构 ${}'\indf^r_X\simeq\indf^r_X\circ\Xi_X$。
\end{othm}

\begin{proof}
由同构 $\indf^r_X\simeq{}'\indf^r_X\circ\Psi_X$ 与 $\Psi_X\circ\Xi_X\simeq\mathrm{Id}_{\QCoh(X)}$ 推出。
\end{proof}

\begin{othm}{推论}{4.6.9}{n:4.6.9}
函子 $({}'\indf^r_X,{}'\oblv^r_X)$ 构成伴随对。
\end{othm}

\begin{proof}
由引理 4.6.8 通过伴随性形式地推出。
\end{proof}

\begin{othm}{注}{4.6.10}{n:4.6.10}
除非 $X$ 最终余连接，函子 $({}'\indf^r_X,{}'\oblv^r_X)$ 不是伴随对。事实上，若 $X$ 不最终余连接，则函子 ${}'\indf^r_X$ 不保持紧对象：它把 $\mathcal{O}_X\in\QCoh(X)$ 送到 $\Crys^r(X)$ 的非紧对象。
\end{othm}

\begin{proof}[命题 4.6.4 的替代证明]
由推论 4.6.9，命题 4.6.4 的断言（在最终余连接情形）等价于：函子 ${}'\indf^r_X$ 的本质像生成 $\Crys^r(X)$。然而后者由 $\indf^r_X$ 的相应事实显然成立。
\end{proof}

\ob{4.6.11}{n:4.6.11}设 $X$ 是最终余连接的 DG 概形，考虑伴随函子对
\[
\Xi_X:\QCoh(X)\rightleftarrows\IndCoh(X):\Psi_X,
\]
其中 $\Xi_X$ 全忠实（见 [IndCoh, Sect.~1.4]）。

我们已经看到，函子 $\indf^r_X$ 通过余局部化函子 $\Psi_X$ 分解。然而，函子 $\oblv^r_X$ 一般并不通过 $\Xi_X$ 分解，即不取值于 $\QCoh(X)$（经 $\Xi_X$ 看作 $\IndCoh(X)$ 的全子范畴）中。

事实上，有下述结论：

\begin{othm}{引理}{4.6.12}{n:4.6.12}
（Drinfeld）函子 $\oblv^r_X$ 通过 $\QCoh(X)$ 在 $\Xi_X$ 下的本质像分解，当且仅当 $X$ 是 Gorenstein 的。
\end{othm}

回忆：DG 概形 $X$ 称为 Gorenstein 的，如果

(a) $\omega_X\in\Coh(X)$（这等价于 $X$ 最终余连接，见 [IndCoh, Proposition 9.6.11]）；

(b) 看作凝聚层时，$\omega_X$ 是分次线丛（这等价于 $\omega_X\in\QCoh(X)^{\mathrm{perf}}$，见 [IndCoh, Corollary 7.4.3]）。

\begin{annote}
条件 (b) 中的「分次线丛」是导出意义下的说法：$\omega_X$ 在同伦意义下可逆，即属于 $\QCoh(X)^{\mathrm{perf}}$。光滑概形上 $\omega_X$ 是线丛，条件自动满足；因此 Gorenstein 条件刻画的是「对偶化复形不比线丛更坏」的奇异性。\hfill【来源：原书 §4.6.11--4.6.12；[IndCoh, Sect.~9.6]（A）】
\end{annote}


\begin{proof}[引理 4.6.12 的证明]
假设 $\oblv^r_X$ 通过 $\QCoh(X)$ 分解。特别地，得到 $\omega_X\in\Coh(X)$ 落在 $\Xi_X$ 的本质像中。现在断言由 [IndCoh, Lemma 1.5.8] 推出。

对反向蕴含，把 $\oblv^r_X(M)$ 写成
\[
\Upsilon_X(M)=\oblv^l_X(M)\otimes\omega_X,
\]
其中张量积按 $\QCoh(X)$ 对 $\IndCoh(X)$ 的作用理解，见 [IndCoh, Sect.~1.4]。还要回忆：函子 $\Xi_X$ 按定义与上述 $\QCoh(X)$-作用相容。因此，若 $\omega_X$ 作为完美对象落在 $\Xi_X$ 的本质像中，则 $\oblv^l_X(M)\otimes\omega_X$ 也落在其中。
\end{proof}

\section{与 Abel 范畴的关系}

本小节中，设 $X$ 是仿射 DG 概形。我们将把范畴 $\Crys^r(X)$ 与一个更熟悉的对象联系起来。

\ob{4.7.1}{n:4.7.1}由于 $\Crys^r(X)$ 上的 t-结构与滤过余极限相容，得到 $\Crys^r(X)^{\heartsuit}$ 是 Grothendieck Abel 范畴。

利用 $\Crys^r(X)$ 在其 t-结构下右完备这一事实，由 [Lu2, Theorem 1.3.2.2] 中的箭头反转，得到典范定义的 t-正合函子
\begin{equation}
D\big(\Crys^r(X)^{\heartsuit}\big)^{+}\to\Crys^r(X)^{+},\tag{4.5}
\end{equation}
其中 $D(-)^{+}$ 表示 Abel 范畴的导出范畴的最终余连接部分。

\ob{4.7.2}{n:4.7.2}我们将证明：

\begin{othm}{命题}{4.7.3}{n:4.7.3}
函子 (4.5) 唯一地延拓为范畴的等价
\[
D\big(\Crys^r(X)^{\heartsuit}\big)\to\Crys^r(X).
\]
\end{othm}

本小节余下部分专门证明命题 4.7.3。不妨设 $X$ 是经典的。

\ob{4.7.4}{n:4.7.4}第 1 步。先假设 $X$ 是光滑经典概形。此时断言由下述事实显然：
\[
\indf^r_X(\mathcal{O}_X)
\]
是两个范畴的紧投射生成元。

\ob{4.7.5}{n:4.7.5}第 2 步。证明函子
\[
D\big(\Crys^r(X)^{\heartsuit}\big)^{+}\to\Crys^r(X)^{+}
\]
是等价。为此，只需证明每个对象 $M\in\Crys^r(X)^{\heartsuit}$ 可以嵌入一个内射对象，即一个 $I\in\Crys^r(X)^{\heartsuit}$，使得
\[
H^0(N)=0\Rightarrow\Hom_{\Crys^r(X)}(N,I)=0.
\]
取闭浸入 $i:X\hookrightarrow Z$，其中 $Z$ 是光滑经典概形。取嵌入 $i_{\dR,*}(M)\hookrightarrow J$，其中 $J$ 是 $\Crys^r(Z)$ 中的内射对象（按同样的意义）；由第 1 步，它存在。

由于函子 $\indf^r_Z$ t-正合，得到 $\oblv^r_Z(J)$ 是 $\QCoh(Z)^{\heartsuit}$ 的内射对象。这蕴含 $I:=i^{\dagger,r}(J)$ 属于 $\Crys^r(X)^{\heartsuit}$ 并具有所需性质。

\ob{4.7.6}{n:4.7.6}第 3 步。注意由推论 4.5.4，范畴 $\Crys^r(X)$ 与 $\Crys^r(X)^{+}$ 的左完备化等同。因此，只需证明典范嵌入
\[
D\big(\Crys^r(X)^{\heartsuit}\big)^{+}\hookrightarrow D\big(\Crys^r(X)^{\heartsuit}\big)
\]
把 $D(\Crys^r(X)^{\heartsuit})^{+}$ 与 $D(\Crys^r(X)^{\heartsuit})$ 的左完备化等同。

为此，只需给出 $\Crys^r(X)^{\heartsuit}$ 的一个有界 Ext 维数的生成元 $P$。

考虑对象
\[
P:=\indf^r_X(\mathcal{O}_X).
\]
它由命题 4.2.11(b) 具有所需性质。

\begin{othm}{注}{4.7.7}{n:4.7.7}
标准的论证可以把命题 4.7.3 的断言推广到对角仿射的拟紧 DG 概形 $X$。
\end{othm}

\begin{othm}{注}{4.7.8}{n:4.7.8}
一旦我们在光滑仿射经典概形上把晶体与 D-模等同起来，就会得到拟紧 DG 概形上 $\Crys^r(X)$ 的许多其他性质：例如，Abel 范畴 $\Crys^r(X)^{\heartsuit}$ 是局部 Noether 的\footnote{这里的意思是：$\Crys^r(X)^{\heartsuit}$ 由它的紧对象生成，且紧对象的子对象是紧的。}，且 $\Crys^r(X)$ 关于其 t-结构有有限上同调维数\footnote{这里的意思是：存在 $N\in\mathbb{N}$，使得对 $n>N$ 及 $\mathcal{M}_1,\mathcal{M}_2\in\Crys^r(X)^{\heartsuit}$，有 $\Hom_{\Crys^r(X)}(\mathcal{M}_1,\mathcal{M}_2[n])=0$。}。注意由命题 4.2.5，为建立这两条性质，只需对仿射光滑经典概形证明它们。
\end{othm}
